\newpage
\phantomsection
\label{prf:restriction-preserves-injective}
\LRAProofFor{prop:restriction-preserves-injective}

\begin{remark*}[Return]
\hyperref[prop:restriction-preserves-injective]{Return to Proposition}
\end{remark*}

\begin{proposition*}[Domain Restriction Preserves Injectivity]
If $f:A\to B$ is injective and $C\subseteq A$, then $f|_C:C\to B$ is injective.
\end{proposition*}

\begin{proof}[Professional Standard Proof]
\LRAProofBodyStart
Let $c_1,c_2\in C$ and suppose $f|_C(c_1)=f|_C(c_2)$. Then
$f(c_1)=f(c_2)$. Since $f$ is injective, $c_1=c_2$.
\end{proof}

\begin{proof}[Detailed Learning Proof]
\LRAProofBodyStart
A restriction has fewer inputs but the same values on those inputs. If two
surviving inputs collide under the restriction, they also collide under the
original function. Injectivity of the original function rules this out.
\end{proof}
\begin{remark*}[Proof structure]
\end{remark*}



\begin{dependencies}
\begin{itemize}
  \item \hyperref[prop:restriction-preserves-injective]{Proof target}
\end{itemize}
\end{dependencies}

\clearpage

